% Section 5.3, from lectures/L05/README.md: the objectives, the evaluation
% questions, and the next lecture.

\section{Review}
\label{c5:sec:review}

\needspace{6\baselineskip}
\subsection*{What you should be able to do now}
After this chapter, you should:

\begin{itemize}
  \item Understand how Zephyr turns a devicetree node into a \code{struct device} at compile time.
  \item Be able to write a driver that serves every instance of its compatible in the devicetree.
  \item Understand the order devices are initialized in, and how to find it for a build.
  \item Be able to implement a subsystem API, so that applications use the part without knowing it.
  \item Be able to package drivers, bindings and Kconfig as a Zephyr module.
  \item Be able to write an emulator that lets a driver run, and fail, without hardware.
  \item Be able to test a driver with ztest, and run the tests with Twister.
\end{itemize}

\needspace{6\baselineskip}
\subsection*{Questions to test yourself}
You should be able to explain:

\begin{enumerate}
  \item What is the difference between a device’s \code{config} and its \code{data}, and where does each one
    live?
  \item How does \code{DT_DRV_COMPAT} connect a driver to the devicetree nodes it serves?
  \item What decides the order in which devices are initialized, and how can you see it for a build?
  \item What does \code{device_is_ready()} return for a device whose init function failed?
  \item Why does the driver’s Kconfig option depend on \code{DT_HAS_QACADEMY_THERMO_ENABLED}?
  \item Why does a sensor driver split reading into \code{sample_fetch} and \code{channel_get}?
  \item What does an emulator make testable that a real part does not?
  \item What is the difference between a ztest suite and a Twister scenario?
  \item What goes wrong without the module’s \code{vendor-prefixes.txt}, and why only under Twister?
\end{enumerate}

\needspace{6\baselineskip}
\subsection*{Looking ahead}
Chapter~\ref{c6:ch} is about subsystems and debugging:

\begin{itemize}
  \item The shell, with commands of your own.
  \item Settings that survive a restart, on the simulated flash.
  \item Networking through the host’s sockets.
  \item Debugging on the host: GDB, sanitizers, fatal errors and coverage.
\end{itemize}

You will make the program talk back, remember things, and crash usefully.
